\chapter{The Volatility-Index Complex}\label{ch:s2:the-volatility-index-complex}

VIX ended 2015 2.4\% above where it began. A long position in VIX futures, held at a constant maturity of about a month and rebuilt here from
Cboe's daily settlements, lost 37\% over the same year. It lost in twelve of the fourteen calendar years from 2013 to September 2026, 99.99\% in all. The futures
curve slopes up most of the time, and a long position rolls down it: the roll-down pays the short sellers. On 5 February 2018 the short sellers
learned the other side of the trade. Our constant-maturity index rose 95.9\% in a day, and an inverse product on it lost the same, 95.9\%. This chapter
builds the complex on the synthetic variance path of chapter~1: a VIX-like index, its futures, long, inverse and leveraged products, and the futures
those products must trade at every close. The build is \texttt{firm.vixetp}.

\section{VIX futures and their curve}

Cboe designs VIX to measure the market's expectation of 30-day forward-looking volatility of the S\&P 500. It computes it from the prices of SPX and
SPXW options: a strip of out-of-the-money strikes weighted as in a variance swap's replication, at a near and a next term, interpolated to a constant
30 days. A VIX future (Book~1, chapter~25) settles at expiry to a Special Opening Quotation of the index, computed that morning from
SPX options of a single expiry 30 days away. Before then it trades at the market's price of
the index's level at expiry, not today's.

\begin{definition}[VIX term-structure slope]\label{def:s2:the-volatility-index-complex:slope}
The \emph{VIX term-structure slope}\index{VIX term-structure slope} is the difference, or ratio, between a later and an earlier VIX future (often the
second and the first contract) or between the first contract and the index; a positive slope (contango) means that longs pay a roll-down as each
contract converges to the index.
\end{definition}

From January 2013 to September 2026 the second contract settled above the first on 85.0\% of days, by 0.89 points on average, and the first contract
stood 0.63 points above VIX. Our long index is rebuilt at each close
between the first two contracts, with weights in proportion to the days left, so that its maturity stays one contract cycle, about a month.
A long position of constant maturity must keep selling the nearer contract and buying the further one. When the curve is
upward-sloping, each future drifts down towards the index as its expiry nears. That drift is the volatility roll-down (Book~5, chapter~25). In our
reconstruction the long index lost at a log rate of 66.4\% a year, 99.99\% over the whole period.

\begin{omfigure}
\begin{tikzpicture}
\begin{axis}[width=12cm,height=5.6cm,ybar,bar width=4pt,ylabel={\small \% over the year},tick label style={font=\footnotesize},
  xmin=2012.4,xmax=2026.6,ymin=-100,ymax=180,xtick={2013,2015,2017,2019,2021,2023,2025},grid=major,grid style={gray!25},
  table/col sep=comma,/pgf/number format/1000 sep={},legend style={font=\scriptsize,at={(0.5,-0.18)},anchor=north},legend columns=2]
\addplot[fill=omThm!70,draw=omThm] table[x=year,y=long]{figdata/strategies-2/05-the-volatility-index-complex/yearly.csv};
\addplot[fill=omDef!60,draw=omDef] table[x=year,y=vix]{figdata/strategies-2/05-the-volatility-index-complex/yearly.csv};
\legend{constant-maturity long in VIX futures,change in VIX}
\end{axis}
\end{tikzpicture}

\omcaption{Each calendar year from 2013 to 2026 (2026 to 24 September): the return of a long position in VIX futures held at a constant maturity of
one contract cycle, rebuilt from Cboe's daily settlements, and the change in VIX. The long position lost in every year but 2018 and 2020. Derived from Cboe's
futures and index histories; the raw data are not redistributed. Data: \texttt{s2\_fetch\_vx}.}\label{fig:s2:the-volatility-index-complex:yearly}
\end{omfigure}

The index is a year-by-year illustration of what the roll-down does (\cref{fig:s2:the-volatility-index-complex:yearly}). In 2014 VIX rose 35\% and the
long position lost 25\%. In 2025 VIX fell 17\% and the long position lost 45\%. Only in 2018 and 2020, years with a spike large enough to outrun the
roll, did the long position gain.

\subsection*{The synthetic complex}

\texttt{firm.vixetp} builds the same objects on chapter~1's variance path, with a calmer volatility of variance (0.25 instead of 0.5). The index is the
21-day implied vol, the expected variance times the premium, in points. Futures expire every 21 trading days. Each is priced at the index level
expected at expiry, reverting to its mean at the variance's speed, times $1 + 0.5h$ for $h$ years to expiry: a term premium that longs pay
(\cref{lst:s2:the-volatility-index-complex:curve}). The index averaged 19.6, with a median of 17.7. The second contract stood above the first on
81\% of days, and the first above the index by 0.44 points. The long constant-maturity index lost at a log rate of 74\% a year, with a volatility of
88\%. Without the term premium, the curve still slopes up on 65\% of days, because the index spends most of its time below its mean. The long index
then loses at a log rate of 26\% a year, though its average daily return is close to zero ($-3.7\%$ a
year): the loss comes from compounding a series this volatile.

\section{Exchange-traded volatility products and their rebalancing}

\begin{definition}[Exchange-traded volatility product]\label{def:s2:the-volatility-index-complex:etp}
An \emph{exchange-traded volatility product}\index{exchange-traded volatility product} is an exchange-traded fund or note (Book~3, chapter~19) that
delivers a multiple ($1$, $-1$, $2$, $-0.5$) of the daily return of an index of VIX futures, usually of constant maturity; it holds or hedges the
futures and rebalances them at each close.
\end{definition}

A product that promises $L$ times the daily return, with value $A$, holds $LA$ of futures. After a daily return $r$, the value is $A(1 + Lr)$ and the
futures held are worth $LA(1 + r)$. To hold $L$ times the new value it must buy
\[
LA(1 + Lr) - LA(1 + r) = A\,L(L - 1)\,r.
\]

\begin{definition}[ETP rebalancing flow]\label{def:s2:the-volatility-index-complex:flow}
The \emph{ETP rebalancing flow}\index{ETP rebalancing flow} is the quantity of futures that leveraged and inverse volatility products must trade at
the close to restore their multiples, $A\,L(L - 1)\,r$ for each product; it is zero for $L = 1$ and has the sign of the day's move for $L = -1$ and
$L = 2$, so these products buy after rises and sell after falls.
\end{definition}

For an inverse product, $L(L - 1) = 2$. A day on which the futures index rises 50\% forces it to buy futures worth its whole value at the close, and
twice-leveraged longs must do the same. The flow is known in advance, as a formula in the day's move, and it arrives at the close when liquidity is
thinnest. Whaley's warning came early: the most popular VIX products are not suitable buy-and-hold investments and are virtually guaranteed to lose
money over time. Products on the short-term futures indices had lost nearly \$4 billion since their launch in 2009.

\section{Futures against options}

A VIX future and SPX options price related things. The options price the S\&P 500's variance over a window, and the future prices the square root of
a 30-day variance that starts at its expiry. The square root is concave, so a future on volatility is worth less than the square root of the forward
variance: the difference is a convexity adjustment, which grows with the volatility of volatility. A trader who thinks the adjustment in the futures
is wrong can trade the future against a forward-starting variance built from SPX options at two expiries. What they hold is the volatility of
volatility, and VIX options (Book~5, chapter~14) price it directly. When the index spikes, all three move together. The relative-value book's risk is
that they then move by different amounts: the future settles at the Special Opening Quotation, the options at their own.

\section{February 2018}

The complex was tested on 5 February 2018. VIX closed at 17.31 on Friday 2 February and at 37.32 on Monday 5 February. The first future settled at
15.625, then 33.225. Our constant-maturity index rose 95.9\% on the Monday and fell 25.8\% on the Tuesday, its best and worst days of 2013 to 2026. An
inverse product on it had lost 96.7\% since the start of the year.

Credit Suisse announced on 6 February that its inverse VIX note, XIV, had suffered an acceleration event: its intraday indicative value on 5 February
had been at or below 20\% of the previous day's closing indicative value, \$108.3681. The note was to be redeemed at its value on the accelerated
valuation date. Augustin, Cheng and Van den Bergen described the episode as a lesson in hedge and leverage rebalancing when markets are concentrated
and volatile, and compared it with portfolio insurance. On our reconstruction's move of 95.9\%, an inverse product's rebalancing flow is 1.92 times its value, all of it
bought at the close.

The synthetic complex has its own February. After six years of calm the first planted crash takes the index from 17.3 to 63.6 in a day, and the long
constant-maturity index rises 189\%. An inverse product with the acceleration clause had grown to 16.5 times its starting value; it ends that day at
zero. At that close it must buy futures worth 3.77 times its previous value, twice the day's move, and so must a twice-leveraged long
(\cref{fig:s2:the-volatility-index-complex:products}). A short-futures book is the same trade at a chosen size:

\begin{center}
\small
\begin{tabular}{@{}lcccc@{}}
\toprule
short futures, rebalanced daily & annual return & maximum drawdown & final value (start 1) & survives\\
\midrule
0.10 of capital & 4.8\% & $-27\%$ & 2.57 & yes\\
0.25 of capital & 9.6\% & $-67\%$ & 6.26 & yes\\
0.50 of capital & --- & $-100\%$ & 0 & no\\
1.00 of capital & --- & $-100\%$ & 0 & no\\
\bottomrule
\end{tabular}
\end{center}

\begin{omfigure}
\begin{tikzpicture}
\begin{axis}[width=11cm,height=5.6cm,ymode=log,xlabel={\small year},ylabel={\small value (start 1)},tick label style={font=\footnotesize},
  xmin=0,xmax=20,ymin=0.5,ymax=30,grid=major,grid style={gray!25},table/col sep=comma,unbounded coords=jump,
  legend style={font=\scriptsize,at={(0.5,-0.3)},anchor=north},legend columns=2]
\addplot[omThm,thick] table[x=year,y=inverse]{figdata/strategies-2/05-the-volatility-index-complex/products.csv};
\addplot[omDef,thick] table[x=year,y=quarter]{figdata/strategies-2/05-the-volatility-index-complex/products.csv};
\legend{inverse product (ends at the first crash),short futures at 0.25 of capital}
\end{axis}
\end{tikzpicture}

\omcaption{Two short positions in the synthetic constant-maturity VIX futures index, on a log scale: an inverse product with an acceleration clause at
20\% of the previous close, which grows to 16.5 times its value and ends at the first planted crash (year 6.0), and a book short a quarter of its
capital, which survives three crashes. Data: \texttt{s2\_vixetp.paths}.}\label{fig:s2:the-volatility-index-complex:products}
\end{omfigure}

Size is the whole strategy. At a quarter of its capital the book earns 9.6\% a year and survives a 67\% drawdown. At half, it is wiped out. The
roll-down is real, but its magnitude says nothing about the size that survives the day the index triples.

\section{Strategy files}

\begin{strategyfile}{Short VIX futures roll-down}\label{strat:s2:the-volatility-index-complex:rolldown}
\sfield{Who pays you, and why} Buyers of volatility futures, directly or through long products, who pay a term premium for protection that is
there when they need it.
\sfield{Instruments and venues} VIX futures on Cboe Futures Exchange; inverse products.
\sfield{Signal} The slope of the curve: short when in contango, smaller or flat when inverted.
\sfield{Sizing and execution} A small fraction of capital, set by a spike that triples the index; daily rebalancing of the maturity.
\sfield{Costs} Futures spreads; roll costs.
\sfield{How it dies} A spike: 5 February 2018 ended an inverse note whose value fell more than 80\% in a day.
\sfield{Horizon, capacity, infrastructure} Months to years; futures access and a daily roll.
\sfield{Backtest honestly} Include 2018 and 2020; use settlements, not index levels; size on the worst day, not the Sharpe ratio.
\sfield{Sources} Whaley (2013); Credit Suisse (2018); Augustin, Cheng and Van den Bergen (2021); this chapter: a long constant-maturity position
lost 99.99\% from 2013 to 2026, and the inverse lost 96.7\% from 2 January to 5 February 2018.
\end{strategyfile}

\begin{strategyfile}{ETP rebalancing anticipation}\label{strat:s2:the-volatility-index-complex:flow}
\sfield{Who pays you, and why} Leveraged and inverse products whose rebalancing, $A\,L(L - 1)\,r$, is a formula in the day's move and trades at the
close.
\sfield{Instruments and venues} VIX futures before and at the close.
\sfield{Signal} The products' published values and multiples, and the day's move so far.
\sfield{Sizing and execution} Small against the flow; positions taken before the close and unwound after.
\sfield{Costs} Spreads at the close, where liquidity is thin.
\sfield{How it dies} Products shrink or close; others anticipate the same flow. Trading ahead of a known flow is legal when it uses public
information; trading on a client's order is not (chapter~29).
\sfield{Horizon, capacity, infrastructure} Hours; a live estimate of the products' flows.
\sfield{Backtest honestly} Close-price fills are optimistic on the days that matter.
\sfield{Sources} Augustin, Cheng and Van den Bergen (2021) on rebalancing risk; no performance figure verified.
\end{strategyfile}

\begin{strategyfile}{VIX futures against SPX options}\label{strat:s2:the-volatility-index-complex:options}
\sfield{Who pays you, and why} Mispricing of the convexity adjustment between the future and forward variance from SPX options.
\sfield{Instruments and venues} VIX futures; SPX options at two expiries; VIX options for the volatility of volatility.
\sfield{Signal} The future against the square root of the forward variance, adjusted for convexity.
\sfield{Sizing and execution} Vega-matched legs; stressed on a spike.
\sfield{Costs} Two markets' spreads; settlement differences.
\sfield{How it dies} A spike in which the legs move by different amounts; settlement at different prices.
\sfield{Horizon, capacity, infrastructure} Weeks; SPX surfaces and VIX futures.
\sfield{Backtest honestly} Model the Special Opening Quotation at expiry.
\sfield{Sources} No performance figure verified.
\end{strategyfile}

\begin{strategyfile}{Long volatility on an inverted curve}\label{strat:s2:the-volatility-index-complex:inverted}
\sfield{Who pays you, and why} Nobody pays, on average: this is insurance that costs little when the curve is inverted, because the roll then favours
longs.
\sfield{Instruments and venues} VIX futures or calls.
\sfield{Signal} An inverted curve (first contract above the second).
\sfield{Sizing and execution} Small; held while inverted.
\sfield{Costs} Spreads; wider in spikes.
\sfield{How it dies} The curve reverts to contango quickly and the roll turns against the long.
\sfield{Horizon, capacity, infrastructure} Days to weeks.
\sfield{Backtest honestly} Inversions are rare (15\% of days from 2013 to 2026); results rest on a few episodes.
\sfield{Sources} This chapter's curve statistics; no performance figure verified.
\end{strategyfile}

\section{Tutorial: roll-down}\label{tut:s2:the-volatility-index-complex:tut}

\begin{tutorial}
\textbf{Goal.} Build the synthetic index, its futures and products; measure the roll-down, the rebalancing flow and the first crash; rebuild the real
constant-maturity index from Cboe's settlements. \textbf{End state:} the table and the two figures.
\begin{tutsteps}
\item \textbf{Futures and the constant-maturity index}.
\omcode{code/firm/vixetp/firm_vixetp.py}{36}{62}{Futures with a term premium and a long index of constant maturity.}\label{lst:s2:the-volatility-index-complex:curve}
\item \textbf{Products and their flows}.
\omcode{code/firm/vixetp/firm_vixetp.py}{65}{82}{A daily-multiple product with an acceleration clause, and its rebalancing flow.}\label{lst:s2:the-volatility-index-complex:etp}
\item \textbf{Run} \texttt{s2\_fetch\_vx.py} once (about 170 files; set \texttt{OQB\_CACHE}), then \texttt{curve\_stats()}, \texttt{short\_book()},
\texttt{spike()} and \texttt{fig\_vixetp.py}.
\end{tutsteps}
\textbf{What to change next.} Scale the short book by the curve's slope; add a VIX call hedge and price it; give the products assets and the futures
an open interest, and measure the flow as a share of it.
\end{tutorial}

\section{Build: volatility index and products}\label{bld:s2:the-volatility-index-complex:vixetp}

\begin{build}
\bfield{Purpose} A VIX-like index, its futures, a constant-maturity index, daily-multiple products and their rebalancing flows.
\bfield{Interface} \texttt{vix\_index(v, cfg)}, \texttt{future(x, mean, cfg, h, premium)}, \texttt{curve\_path(v, cfg, premium, every)},
\texttt{etp(r, leverage, accel)}, \texttt{flow(product, r, leverage)}.
\bfield{Rules} The expiring leg settles at the index; weights set at the previous close; a product ends when a day leaves it at or below the
acceleration fraction of the previous close.
\bfield{Acceptance tests} \texttt{code/firm/vixetp/tests/}: the flow formula by hand for $L = -1, 1, 2$; an ended product frozen with no further
flows; a flat index with no premium gives a flat curve and no roll, and with a premium a contango and a steady roll-down.
\bfield{Stretch} A second, faster volatility factor; VIX options on the index; products' assets and open interest.
\end{build}

\begin{omsources}
\item Cboe, Volatility Index Methodology: Cboe Volatility Index.
\item R.~E.~Whaley, ``Trading volatility: at what cost?'', \emph{Journal of Portfolio Management} 40(1), 2013.
\item P.~Augustin, I.-H.~Cheng and L.~Van den Bergen, ``Volmageddon and the failure of short volatility products'', \emph{Financial Analysts Journal}
77(3), 2021.
\item Credit Suisse AG, ``Credit Suisse AG announces the event acceleration of its XIV ETNs'', media release, 6 February 2018 (SEC exhibit 99.1).
\item Cboe Futures Exchange, VIX futures daily settlements; Cboe, VIX history.
\end{omsources}

\section{Exercises}

\begin{exercise}[$\star$]\label{exo:s2:the-volatility-index-complex:1}
A product returns $-1$ times the index each day and is worth 100. The index rises 30\%. What must it buy at the close?
\end{exercise}

\begin{exercise}[$\star$]\label{exo:s2:the-volatility-index-complex:2}
The same for a product returning twice the index, and for one returning the index once.
\end{exercise}

\begin{exercise}[$\star$]\label{exo:s2:the-volatility-index-complex:3}
The first contract has 10 days to expiry and the second 38, with contracts 28 days apart. What weight does a constant 30-day position put on the first?
\end{exercise}

\begin{exercise}[$\star\star$]\label{exo:s2:the-volatility-index-complex:4}
In 2014 VIX rose 35\% and the constant-maturity long lost 25\%. Explain.
\end{exercise}

\begin{exercise}[$\star\star$]\label{exo:s2:the-volatility-index-complex:5}
Why is a future on volatility worth less than the square root of the forward variance?
\end{exercise}

\begin{exercise}[$\star\star$]\label{exo:s2:the-volatility-index-complex:6}
Why do inverse and leveraged volatility products make a spike worse at the close?
\end{exercise}

\begin{exercise}[$\star\star\star$]\label{exo:s2:the-volatility-index-complex:7}
\emph{Coding.} Run \texttt{curve\_stats(0.0)}. Without a term premium, why does the curve still slope up on most days, and why does the long index
still lose?
\end{exercise}

\begin{exercise}[$\star\star\star$]\label{exo:s2:the-volatility-index-complex:8}
\emph{Find the flaw.} ``Short VIX futures made 9.6\% a year at a quarter of capital, so at the full capital it would make almost 40\%.''
\end{exercise}

\section{Problem: Roll-Down}

\begin{problem}[{Weekend problem --- short volatility through futures}]\label{pb:s2:the-volatility-index-complex:1}
The synthetic complex, Cboe's futures and February 2018.

\medskip\noindent\textbf{Part I --- The curve.}
\begin{enumerate}
\item What does VIX measure, and how is it computed?
\item Define the \omterm{def:s2:the-volatility-index-complex:slope}{VIX term-structure slope} and give its real statistics.
\item What is the roll-down, and what did it cost a constant-maturity long?
\item What did the long lose in 2014, 2015 and 2025, and what did VIX do?
\end{enumerate}

\medskip\noindent\textbf{Part II --- Products.}
\begin{enumerate}[resume]
\item Define an \omterm{def:s2:the-volatility-index-complex:etp}{exchange-traded volatility product}.
\item Derive the rebalancing flow $A\,L(L - 1)\,r$.
\item Which products buy after rises?
\item What did Whaley find?
\end{enumerate}

\medskip\noindent\textbf{Part III --- February 2018.}
\begin{enumerate}[resume]
\item What happened to VIX and the first future on 5 February?
\item What did the constant-maturity index and its inverse do?
\item What did Credit Suisse announce, and why?
\item How did Augustin, Cheng and Van den Bergen read the episode?
\end{enumerate}

\medskip\noindent\textbf{Part IV --- The verdict.}
\begin{enumerate}[resume]
\item State the \emph{named result}: the short-futures book's roll-down return and the inverse product's loss in the planted spike.
\item What did the short books earn at each size?
\item What flow did the inverse product face at the crash's close?
\item Why is the term premium not needed for the curve to slope up?
\item How would you size a short VIX futures book?
\item How would you backtest it honestly?
\item Which strategy file profits from the others' rebalancing?
\item In one sentence: who pays the roll-down, and why?
\end{enumerate}
\end{problem}

\section{Interview questions}

\begin{interviewq}[$\star$ \iqroles{trader}]\label{iq:s2:the-volatility-index-complex:1}
Why do long VIX products lose money over time?
\end{interviewq}

\begin{interviewq}[$\star\star$ \iqroles{trader}]\label{iq:s2:the-volatility-index-complex:2}
The VIX curve is inverted. What does it tell you, and what trades does it suggest?
\end{interviewq}

\begin{interviewq}[$\star\star$ \iqroles{researcher}]\label{iq:s2:the-volatility-index-complex:3}
How would you build a constant-maturity index from futures settlements, and what can go wrong at the roll?
\end{interviewq}

\begin{interviewq}[$\star\star$ \iqroles{risk}]\label{iq:s2:the-volatility-index-complex:4}
You run a short VIX futures book. How do you size and stress it?
\end{interviewq}

\begin{interviewq}[$\star\star$ \iqroles{developer}]\label{iq:s2:the-volatility-index-complex:5}
Design the daily process that computes a leveraged product's rebalancing at the close.
\end{interviewq}

\begin{interviewq}[$\star\star\star$ \iqroles{researcher}]\label{iq:s2:the-volatility-index-complex:6}
Show that a product returning $L$ times the daily return must trade $A\,L(L - 1)\,r$ at each close, and show that its value over two days is not $L$
times the index's two-day return.
\end{interviewq}
